\section*{6 Query-efficient distance amplification}
In this section we construct our query-efficient distance amplification procedure. We start by giving a somewhat more formal definition of locally decodable codes (compared to Definition 1) or, more precisely, a more formal definition of their non-adaptive counterparts. Recall that, informally, these are LDC in which the joint distribution of queries depends solely on the index one wishes to decode and is independent of the received word. By inspection, it is our understanding that the AEL distance amplification procedure also requires non-adaptivity.
\setcounter{nttheorem}{63}
\begin{ntdefinition}[Author Definition 64: Locally decodable codes]
Let $(C,Q,R)$ be a tuple of functions
\begin{align*}C&:\Sigma_{\mathrm{in}}^k\to\Sigma_{\mathrm{out}}^n,\\Q&:[k]\times\{0,1\}^r\to[n]^q,\\R&:[k]\times\Sigma_{\mathrm{out}}^q\times\{0,1\}^r\to\Sigma_{\mathrm{in}}.\end{align*}
Define $D:[k]\times\Sigma_{\mathrm{out}}^n\times\{0,1\}^r\to\Sigma_{\mathrm{in}}$ as follows. For $v\in[k]$, $y\in\Sigma_{\mathrm{out}}^n$, and $s\in\{0,1\}^r$, let
\[Q(v,s)=(u_1,\ldots,u_q),\qquad D(v,y,s)=R(v,y_{u_1},\ldots,y_{u_q},s).\]
The tuple $(C,Q,R)$ is called a $(q,\delta,\varepsilon)$-locally decodable code (or $(q,\delta,\varepsilon)$-LDC for short) if the following holds. For every $v\in[k]$, $x\in\Sigma_{\mathrm{in}}^k$, and $y\in\Sigma_{\mathrm{out}}^n$ for which $\operatorname{dist}(y,C(x))\leq\delta$, it holds that
\[\Pr_{s\sim U_r}[D(v,y,s)=x_v]\geq1-\varepsilon.\]
We call the function $C$ the encoding function, $Q$ the querying function, and $R$ the reconstruction function. The induced function $D$ is called the decoding function. The parameters $k,n$ are referred to as the message length and the block length, respectively. The sets $\Sigma_{\mathrm{in}},\Sigma_{\mathrm{out}}$ are called the input alphabet and output alphabet, respectively. We will be interested mostly in locally decodable codes in which $\Sigma_{\mathrm{in}}=\Sigma_{\mathrm{out}}$ in which case we refer to both as the alphabet of the code. The parameter $r$ is called the randomness complexity of the LDC. We say the LDC is explicit if all three functions $C,Q,R$ are polynomial-time computable. Note that then the decoding function $D$ is also polynomial-time computable.
\end{ntdefinition}
\subsection*{6.1 The distance amplification procedure}
In this section we present our query-efficient distance amplification procedure. We start by describing the building blocks we use and specify their parameters.
\paragraph{Building blocks}
\begin{itemize}
\item For $i=1,2$ let $(C_i,Q_i,R_i)$ be a $(q_i,\delta_i,\varepsilon_i)$-LDC with message length $k_i$ and block length $n_i$ over the same alphabet $\Sigma$. We denote the rate $k_i/n_i$ of $C_i$ by $\rho_i$.
\item Let $(C_3,Q_3,R_3)$ be a family of $(q_3(k_3),\delta_3(k_3),\varepsilon_3(k_3))$-LDC having rate $\rho_3(k_3)$ for message length $k_3$. The code $C_3$ is also over the alphabet $\Sigma$. We will always work with functions $q_3,\delta_3,\varepsilon_3,\rho_3$ that are monotone. More precisely, $q_3$ and $\rho_3$ are non-decreasing and $\delta_3,\varepsilon_3$ are non-increasing. We sometimes write $q_3,\delta_3,\varepsilon_3,\rho_3$ without mentioning explicitly the message length, and by that refer to the largest $k_3$ considered in the construction for $q_3,\delta_3$ and the smallest $k_3$ when considering $\varepsilon_3,\rho_3$. In any case, we assume (mostly for simplicity) that $\rho_3(k_3)\geq1/2$ for all $k_3$.
\item Set $\ell=n_1/k_2$. Let $G=(L,R,E)$ be a $(\delta_2/2,\delta_1)$-sampler with $|L|=\ell$ and $|R|=r$. Assume $G$ is left-regular with left-degree $d=n_2$. Assume further that every right-vertex $v$ of $G$ has degree $\deg(v)\in[D/2,2D]$, where $D$ is the average right degree $D=\ell d/r=n_1/(r\rho_2)$.
\end{itemize}
\paragraph{How to think of the parameters?}
We think of $C_1$ as the code whose distance $\delta_1$ we wish to amplify. Typically, the code $C_2$ has a much shorter message length $n_2\ll n_1$. In all applications in this paper we take $\delta_2$ to be either constant or slightly sub constant in $n_1$. The code $C_3$ has a larger block length than $C_2$ and, depending on the application, it has either a somewhat smaller or much smaller message length than $n_1$. We typically take $\delta_3\approx\delta_2$. The rates of all three codes is taken to be constant and even close to one. Note that we take $C_3$ to be a family of codes, whereas $C_1$ and $C_2$ are codes with predetermined message lengths. The reason is that the sampler $G$ is not necessarily right-regular, and in the construction, we associate codes from $C_3$ with the right vertices of $G$. Recall, though that the ratio of largest to smallest right-degree is bounded by $4$, so that is a minor technicality.
To describe the LDC that is composed of these building blocks, we need to specify the encoding function, querying function and reconstruction function. We start by describing the encoding function.
\paragraph{The encoding function}
Let $n=\sum_{v\in R}n_v$ where $n_v$ is the block length of the code from the family $C_3$ having message length $k_v=\deg(v)$. We define the function $C:\Sigma^{k_1}\to\Sigma^n$ as follows. Let $x\in\Sigma^{k_1}$.
\begin{enumerate}
\item Compute $y=C_1(x)\in\Sigma^{n_1}$.
\item Partition $y$ to $y=y^{(1)}\circ\cdots\circ y^{(\ell)}$ consecutive blocks, each of length $k_2$. Recall that, indeed, $n_1=\ell k_2$.
\item For every $u\in[\ell]$ compute $z^{(u)}=C_2(y^{(u)})\in\Sigma^{n_2}$.
\item For every $v\in[r]$ and $j\in[\deg(v)]$ let $(u,j')=\Gamma(v,j)\in[\ell]\times[n_2]$. Define the string $w^{(v)}\in\Sigma^{\deg(v)}=\Sigma^{k_v}$ as follows: for $j\in[\deg(v)]$, $(w^{(v)})_j=(z^{(u)})_{j'}$.
\item For every $v\in[r]$ compute $t^{(v)}=C_3(w^{(v)})\in\Sigma^{n_v}$.
\item The output of the encoding function on input $x$ is then defined by $C(x)=t^{(1)}\circ\cdots\circ t^{(r)}\in\Sigma^n$, where as usual we identify $R$ with $[r]$.
\end{enumerate}
By the construction of the encoding function, the message length and block length of the resulted code are $k_1$ and $n$, respectively. From here on we denote $k=k_1$.
\paragraph{The querying function}
We denote the randomness complexity of $C_1,C_2,C_3$ by $r_1,r_2,r_3$, respectively. The randomness complexity of the resulting querying function will be $r=r_1+r_2+r_3$, and the query complexity will be $q\leq q_1q_2q_3$, where $q_3$ is taken to be the maximum query complexity taken over all right vertices. We turn to define the querying function $Q:[k]\times\{0,1\}^r\to[n]^q$ as follows. On inputs $p\in[k]$, $s\in\{0,1\}^r$ we proceed as follows.
\begin{enumerate}
\item Partition $s=s_1\circ s_2\circ s_3$ where $|s_1|=r_1$, $|s_2|=r_2$, $|s_3|=r_3$.
\item Compute $(a_1,\ldots,a_{q_1})=Q_1(p,s_1)\in[n_1]^{q_1}$.
\item For $i=1,\ldots,q_1$
\begin{enumerate}
\item Set $u_i=\lceil a_i/k_2\rceil$ and $b_i=1+((a_i-1)\bmod k_2)$. Informally, $u_i$ is the ``bucket'' in which $a_i$ resides and $b_i$ is its location within the bucket. Note that we start the counting from $1$ rather than $0$, hence the slightly annoying addition and subtraction by $1$ in the definition of $b_i$.
\item Compute $(t_1^{(i)},\ldots,t_{q_2}^{(i)})=Q_2(b_i,s_2)\in[n_2]^{q_2}$.
\item For $j=1,\ldots,q_2$
\begin{enumerate}
\item Let $(v^{(i,j)},\widehat t_j^{(i)})=\Gamma(u_i,t_j^{(i)})\in[r]\times[k_{v^{(i,j)}}]$.
\item Compute $(e_1^{(i,j)},\ldots,e_{q_3}^{(i,j)})=Q_3(\widehat t_j^{(i)},s_3)\in[n_{v^{(i,j)}}]^{q_3}$.
\item As before, we endow the right vertices of the sampler in a fixed (arbitrary) order by identifying $R$ with $[r]$. For $h=1,\ldots,q_3$ set $c^{(i,j,h)}$ to be the absolute location of $e_h^{(i,j)}$ in the ordering of $R$. That is, $c^{(i,j,h)}=e_h^{(i,j)}+\sum_{v<v^{(i,j)}}n_v$.
\end{enumerate}
\end{enumerate}
\item The result is then given by $Q(p,s)=(c^{(i,j,h)})_{(i,j,h)\in[q_1]\times[q_2]\times[q_3]}$.
\end{enumerate}
Note that, indeed, the query complexity $q$ of the querying function defined above is at most $q_1q_2q_3$ where, recall, $q_3=q_3(2D)$. From here on we identify $[q]$ with $[q_1]\times[q_2]\times[q_3]$.
\paragraph{The reconstruction procedure}
We define the reconstruction procedure $R:[k]\times\Sigma^q\times\{0,1\}^r\to\Sigma$ as follows. On inputs $p\in[k]$, $\sigma=(\sigma^{(i,j,h)})_{(i,j,h)\in[q_1]\times[q_2]\times[q_3]}\in\Sigma^q$, and $s\in\{0,1\}^r$, we proceed as follows.
\begin{enumerate}
\item Partition $s=s_1\circ s_2\circ s_3$ where $|s_1|=r_1$, $|s_2|=r_2$, $|s_3|=r_3$ as in the querying function.
\item For $i=1,\ldots,q_1$
\begin{enumerate}
\item For $j=1,\ldots,q_2$
\begin{enumerate}
\item Denote $(z_1,\ldots,z_{q_3})=(\sigma^{(i,j,1)},\ldots,\sigma^{(i,j,q_3)})$.
\item Compute $y_j^{(i)}=R_3(\widehat t_j^{(i)},z_1,\ldots,z_{q_3},s_3)$, where $\widehat t_j^{(i)}=\widehat t_j^{(i)}(p,s)$ as defined in the querying function.
\end{enumerate}
\item Set $x_i=R_2(b_i,y_1^{(i)},\ldots,y_{q_2}^{(i)},s_2)$ where $b_i=b_i(p,s)$ as defined in the querying function.
\end{enumerate}
\item The output is then given by $R(p,\sigma,s)=R_1(p,x_1,\ldots,x_{q_1},s_1)$.
\end{enumerate}
